\section*{Introduction}
The ATP's Under Pressure Rating sums four percentages: break points saved and converted, tiebreaks won and deciding sets won. Many believe some players rise to big moments, yet the rating rarely shows who they are. We ask why: does the official rating measure clutch skill at all, and can a context-adjusted measure (UPR+) find clutch skill that is real, stable and predictive of future results, including among players ranked 51--150?

\section*{Methods}
We used 7,518 Match Charting Project matches (1.28 million points) with Jeff Sackmann's ATP match, ranking and player files. A Markov model of the score gave every point a leverage: how much winning it changes the probability of winning the match. UPR+ is a player's leverage-weighted point performance relative to his own serve and return level in the same match, after removing tour-wide score effects, shrunk for sample size. Reliability was tested with odd/even-match splits against outcomes simulated without player clutch. For players with at least 10 charted matches in the three seasons before year-end 2014--2024 (568 player-seasons), we predicted next-season wins above a points-based Markov expectation (all tour, Challenger and qualifying matches) and next-year ranking change, adjusting for share of points won, ranking, age, charted volume and season, and compared UPR+ with the official rating from tour-level matches.

\section*{Results}
Big points decide matches: outcomes beyond what points won predict track the 15\% highest-leverage points ($r = 0.82$) and barely the rest ($r = -0.14$); 5.3\% of matches are won with fewer points. The official rating cannot isolate this skill (Table~\ref{tab:uprplus}): tiebreak and deciding-set percentages need 119--612 matches to reach reliability 0.5, and the rating tracks ordinary skill ($r = 0.53$ with points won). UPR+ shows that clutch skill is real but small. Between-player variance is 1.7 times chance (simulation maximum 1.4) and odd and even matches agree ($r = 0.26$; simulation maximum 0.20), implying a true spread of 0.43 percentage points per weighted point, or 3.6 points of match-win probability per match. UPR+ predicted the next season: each SD added 1.0 percentage point of wins per match beyond points (95\% CI 0.6--1.4), twice the official rating's 0.5 ($-0.1$ to 1.0), and the top quintile out-performed the bottom by 3.8 points (Figure~\ref{fig:quint}), about 1.8 wins a season. The effect held after adjusting for past wins above expectation. It did not translate into ranking gains ($-4\%$ per SD; $-12\%$ to $+4\%$), including for players ranked 51--150.

\section*{Conclusion}
Clutch skill exists in men's tennis, but it lives in leverage-weighted points, not in the four official ratios. UPR+ measures it with less than half the matches the official rating needs and predicts who will keep beating their points. Rankings appear to already reflect these wins, so clutch sustains results rather than revealing hidden risers.
